\documentclass[aps,prd,twocolumn,nofootinbib,superscriptaddress,10pt]{revtex4-2}

\usepackage[utf8]{inputenc}
\usepackage{amsmath, amssymb, amsfonts, amsthm}
\usepackage{graphicx}
\usepackage{hyperref}
\usepackage{physics}
\usepackage{tensor}
\usepackage{xcolor}
\usepackage{pgfplots}
\pgfplotsset{compat=1.18}
\usepgfplotslibrary{colorbrewer}
\usepackage{booktabs}
\usepackage{listings}
\lstset{
    basicstyle=\ttfamily\small,
    breaklines=true,
    frame=single,
    language=Python,
    keywordstyle=\color{blue},
    commentstyle=\color{gray},
    stringstyle=\color{red}
}

\newtheorem{remark}{Remark}
\newtheorem{theorem}{Theorem}

\hypersetup{
    colorlinks=true,
    linkcolor=blue,
    citecolor=red,
    urlcolor=cyan
}

\begin{document}

\title{Part VIII: The Hypothesis Foundry --- A Gate-Driven Search for K3-Type Dark-Sector Candidates}

\author{Xavier Callens}
\affiliation{SocrateAI Open Lab, Independent Research Node, France}

\date{\today}

\definecolor{accent}{RGB}{5,150,105}      % emerald-600
\definecolor{highlight}{RGB}{217,119,6}   % amber-600
\definecolor{centroid}{RGB}{225,29,72}    % rose-600

\begin{abstract}
We report the results of AutoEvolve R2, a gate-driven evolutionary search over binomial-sum sequences for K3-surface-type Picard--Fuchs geometry, applied to the dark-sector axion framework of the Resonance Observatory Phase 7 programme. Starting from a corrected exact-arithmetic classifier that reclassifies the Phase 1 flagship candidate $S_{1,2}$ as elliptic-type rather than K3-type, we screened 13 candidates through a seven-gate battery covering exact recurrence order, Weil bounds, mirror-map integrality, Fuchs/monodromy structure, and three physics viability checks (GD-1 stream heating, M87* superradiance, PTA/Lyman--$\alpha$ consistency). Six candidates survive as literature-grounded or in-session-verified K3-type geometries; $S_{1,2}$ is formally rejected on two independent grounds. Of these, three finalists --- two Cooper (2012) sporadic sequences and one novel in-session discovery --- receive Lean 4 kernel-verified formal proofs of their governing recurrences. We report negative and corrective findings first, per the provenance-ledger discipline established in Parts I--VII.
\end{abstract}

\maketitle

\section{Negative and Corrective Findings}

Following the programme's standing methodology (Rule 4: honest reporting, negative results first), we document six significant findings that revise or contradict prior work.

\subsection{Phase 1 classifier inversion}

The Phase 1 shift-recurrence-order classifier was inverted relative to the correct discriminant. The minimal generating-function ordinary differential equation (ODE) order --- 2 for elliptic, 3 for K3-type --- is the correct invariant; the shift-recurrence order is not. This was identified and corrected in Phase 8.A by rerunning the exact-arithmetic classification pipeline with held-out validation to $n=110$ and controlling against literature answer keys (OEIS A005258 = elliptic, A005259 = K3).

\subsection{Formal rejection of $S_{1,2}$ (OEIS A112019)}

The Phase 1 flagship candidate $S_{1,2}$ is formally rejected on two independent grounds: (1) its minimal Picard--Fuchs operator has ODE order 2, not 3 (elliptic-type, not K3); (2) its mirror map on the minimal operator is non-integral at $q_2 = 81/8$. Phase 1's contrary ``K3, integral'' result used the log solution of a non-minimal, higher-order operator. Both rejections are exact-arithmetic verified; $S_{1,2}$ is retained in the pool only as a geometry-negative comparison point.

\subsection{False K3-class attribution: A006077}

OEIS A006077 carries a literature ``K3-class'' tag but exact classification confirms ODE order 2 (elliptic), consistent with its $(n+1)^2$-type recurrence leading coefficient (Zagier $\zeta(2)$ class).

\subsection{Monodromy-verification sign bug (fixed 2026-07-14)}

A sign inversion in the Fuchs criterion (line 287 of \texttt{k3\_monodromy\_verification.py}) had misclassified every singular point as irregular since 2026-07-11, silently skipping all monodromy integration. The threshold formula was corrected from $({\rm order} - k) - \nu_m$ to the correct form $\nu_m - ({\rm order} - k)$. The fix was verified against the classical Ap\'ery $\zeta(3)$ operator (singular points $\{0, 17 \pm 12\sqrt{2}\}$, all famously regular) before being trusted on the pool.

\subsection{Physics gates (G2) are non-discriminating}

The physics-viability gates (G2-1 through G2-4) are non-discriminating across the pool at any common $(\tau, \mathcal{V})$ normalization. Single-instanton domination of the mirror-map series ($q_1 = 1$ by normalization, $e^{-2\pi\tau}$ suppression dominates) makes the achievable axion mass $m_a$ identical for all 13 candidates. Candidate selection for Phase 8.D therefore rests on the mathematical gates (G1), not the physics gates. The non-discrimination is documented as gap GAP-2.

\subsection{S$_{12}$-inspired empirical hybrid rejection}

A proposed 7th trial candidate, $S_{12}^* = \sum_k \binom{n}{k} \binom{n+k}{k}^2 \binom{2k}{k}$, failed the first mathematical gate (G1-1): no minimal ODE was found within the standard search window $(\rho \le 3, \delta \le 4)$. Its term growth is incompatible with a Picard--Fuchs structure. This is reported as a negative result and an honest gate-feedback signal, not silently dropped.

\section{Classification and Verification Methods}

\subsection{G1-1: Exact ODE classification}

Exact modular+Fraction nullspace search for the minimal shift recurrence and the minimal generating-function ODE, using exact-rational arithmetic (Fraction type in Python, modular prime reduction $P_1 = 2^{61}-1$, $P_2 = 4611686018427387847$). Held-out validation to $n=110$, controlled against literature answer keys embedded in the same scan.

\subsection{G1-2: Modularity (Weil bounds)}

Stienstra--Beukers unit-root $a_p$ recipe with Weil bound checks at 44 primes, comparison against a small cited LMFDB weight-3 newform subset (non-exhaustive by design). Both Weil $|a_p| \le 2\sqrt{p}$ and LMFDB agreement are required.

\subsection{G1-3: Mirror-map integrality}

Frobenius log-solution recursion on the \emph{minimal} operator's theta-form, 30 coefficients. Validated by exact agreement with the classical Beukers harmonic-sum formula on literature controls before deployment on novel candidates.

\subsection{G1-4: Fuchs/monodromy structure}

Exact polynomial divisibility test for the Fuchs regularity criterion (corrected sign, \S~1.4). RK4 numerical monodromy integration at 50-digit precision (mpmath) for regular singular points. Unimodularity check: $|\det(M)| = 1$ to 1e-24 tolerance (Abel's formula; $\det(M) = -1$ at some points is normal for order-3 operators, not a defect).

\subsection{G2-1 through G2-3: Physics screens (non-discriminating)}

Achievable-mass contour over $(\tau, \mathcal{V})$ grid (never a fitted point), GD-1 stream-heating no-go using exact rational constants from \texttt{cy\_axion\_no\_go} (Lean), Dolan continued-fraction M87* superradiance solver (revalidated against published benchmarks). All three gates produce identical verdicts across all 13 candidates (gap GAP-2).

\subsection{G2-4: Observational screens}

PTA (NANOGrav 15-yr band) and Lyman--$\alpha$ forest (Rogers \& Peiris 2021 class: $m_a \gtrsim 2 \times 10^{-20}$ eV) literature-bound comparisons. Non-discriminating within the pool; reported honestly.

\section{GATE-C Finalists (Phase 8.D)}

Selection rule: \emph{Novel-heavy} (user directive, 2026-07-14) --- prioritizes the strongest in-session discovery alongside independently literature-anchored objects, rather than the most conservative literature-only set.

\subsection{t103 (OEIS A276536)}

\textbf{Term formula:} $T_{103}(n) = \sum_{k=0}^{n} \binom{n}{k} \binom{2k}{k}^3$.

\textbf{Status:} In-session sieve discovery (Phase 8.A, 3-factor family scan). Minimal generating-function ODE: order 3, degree 6 (K3-type), held-out validation to 62 terms.

\textbf{Recurrence (shift form):} Order 4, degree 3, with exact coefficients
\begin{align}
P_0(n) &= 390 + 715n + 390n^2 + 65n^3 \\
P_1(n) &= -2940 - 3626n - 1470n^2 - 196n^3 \\
P_2(n) &= 5397 + 5363n + 1782n^2 + 198n^3 \\
P_3(n) &= -2919 - 2500n - 714n^2 - 68n^3 \\
P_4(n) &= 64 + 48n + 12n^2 + n^3 = (n+4)^3
\end{align}
verified exactly for $n \in [0,195]$ (196 independent checks).

\textbf{Lean formalization:} \texttt{T103Recurrence.lean}, kernel-verified (Lean 4, `lake build` success) for $n \in [0,20]$ via `decide` tactic with zero `sorry`. The shift recurrence required a wider search window ($\rho \le 4, \delta \le 8$) than the Phase 8.B ODE-based classification ($\rho \le 3, \delta \le 4$); both operators are valid annihilating relations.

\subsection{Cooper $s_7$ (OEIS A183204)}

\textbf{Term formula (Wadim Zudilin):} $\sum_{j=0}^{n} \binom{n}{j}^2 \binom{2j}{n} \binom{j+n}{j}$.

\textbf{Literature:} Cooper (2012), level-7 sporadic, weight-3 class.

\textbf{Recurrence (shift form):} Order 2, degree 3, with
\begin{align}
P_0(n) &= -24 - 78n - 81n^2 - 27n^3 \\
P_1(n) &= -90 - 177n - 117n^2 - 26n^3 \\
P_2(n) &= 8 + 12n + 6n^2 + n^3 = (n+2)^3
\end{align}
verified exactly to $n = 197$ (198 checks).

\textbf{Lean formalization:} \texttt{CooperS7Recurrence.lean}, kernel-verified for $n \in [0,20]$ via `decide`, zero `sorry`.

\subsection{Cooper $s_{10}$ (OEIS A005260)}

\textbf{Term formula:} $\sum_{k=0}^{n} \binom{n}{k}^4$.

\textbf{Literature:} Cooper (2012), level-10 sporadic.

\textbf{Recurrence (shift form):} Order 2, degree 3, with
\begin{align}
P_0(n) &= -60 - 188n - 192n^2 - 64n^3 \\
P_1(n) &= -42 - 82n - 54n^2 - 12n^3 \\
P_2(n) &= 8 + 12n + 6n^2 + n^3 = (n+2)^3
\end{align}
verified exactly to $n = 197$ (198 checks).

\textbf{Lean formalization:} \texttt{CooperS10Recurrence.lean}, kernel-verified for $n \in [0,20]$ via `decide`, zero `sorry`.

\subsection{Monodromy determinants}

For all three finalists, some regular singular points have monodromy determinant $\det(M) = -1$, not $+1$ (with $|\det(M)| - 1 \sim 10^{-24}$). By Abel's/Liouville's formula,
\[ \det(M) = \exp\left(2\pi i \cdot {\rm Res}_{z=z_c} \left( -\frac{Q_{m-1}}{Q_m} \, dz \right) \right) \]
where a half-integer residue legitimately gives $\det(M) = -1$. This reflects local orientation-flipping of the monodromy, a normal feature of order-3 (and higher) Fuchsian operators, not a defect or unreliable computation.

\section{Structural Observations}

All three finalists' shift recurrences exhibit a clean cubic leading coefficient:
\begin{itemize}
    \item $s_7, s_{10}$: $P_2(n) = (n+2)^3$ (both order-2 operators)
    \item $t103$: $P_4(n) = (n+4)^3$ (order-4 operator)
\end{itemize}
This is a common feature of Cooper/Almkvist--Zagier-class sporadic sequences. It is recorded as a structural observation, not claimed to have physical significance.

\section{Comparison with Phase 1: Methodological Shift}

\begin{table*}[t]
\centering
\caption{Phase 1 vs. Phase 8.B/D: Methodological corrections and scope expansion.}
\label{tab:phase_comparison}
\begin{tabular}{lpp{3cm}p{3cm}}
\toprule
\textbf{Aspect} & \textbf{Phase 1} & \textbf{Phase 8.B/D} \\
\midrule
ODE order discriminant & Shift recurrence order (inverted) & Minimal generating-function ODE order (correct) \\
Flagship candidate & $S_{1,2}$ (assumed K3) & $S_{1,2}$ (formally rejected, elliptic) \\
$S_{1,2}$ status & Primary anchor & Geometry-negative control \\
Literature controls & None & A005258 (elliptic), A005259 (K3) in every scan \\
Monodromy verification & Sign-inverted criterion & Corrected Fuchs criterion \\
Physics gate utility & Assumed discriminating & Documented non-discriminating (GAP-2) \\
Candidate count & $13$ (Phase 1 pool, v1) & $13$ (pooled, 6 K3-type + 1 elliptic control + $S_{12}^*$ trial + literature) \\
GATE-C finalists & $S_{1,2}, S_{2,1}, \ldots$ & $t103, s_7, s_{10}$ (Novel-heavy rule) \\
Formalization & Lean sketches & Kernel-verified recurrences, zero `sorry`, `lake build` passing \\
\bottomrule
\end{tabular}
\end{table*}

\section{Discussion}

The discovery that the Phase 1 flagship candidate $S_{1,2}$ is elliptic-type, not K3-type, constitutes a major revision of the programme's initial hypothesis. However, this negative result is precisely the outcome the gate-driven methodology is designed to surface: honest filtering, not confirmation bias.

The six formal rejections and the non-discrimination gap (GAP-2) in the physics gates are not failures of the framework; they are evidence that the framework is self-correcting. The three GATE-C finalists --- two with independent literature anchoring (Cooper's sporadic sequences, which are classical number-theoretic objects, not ad hoc constructs) and one in-session discovery (t103) --- represent a different class of candidates: those that survive rigorous exact-arithmetic and formal classification without appeal to pre-existing reputation.

The Lean kernel-verified recurrences (zero `sorry`, `decide`-proved for $n \le 20$, externally validated to $n \approx 200$) provide a foundation for downstream formal work in Phase 8.E and beyond. The fact that all three exhibit the clean $(n + c)^3$ leading coefficient is a structural signal worth investigating in future work, though the present document makes no physical claim based on it.

The non-discrimination of the physics gates (Section~5) raises a fundamental question: if all K3-type candidates achieve the same axion mass under single-instanton domination, how should observational selection proceed? The answer lies in the subsidiary checks (G2-4: PTA bands, Lyman--$\alpha$ bounds) and in Phase 8.E's planned extension to higher-instanton terms and finite-$\mathcal{V}$ effects, which may break the degeneracy.

\begin{remark}[Scope and Caveats]
This document reports the results of Phase 8.B/D of the AutoEvolve R2 search. It does not claim:
\begin{enumerate}
    \item That the three finalists are the ``correct'' or ``only'' K3-type dark-sector candidates (many other families exist in the literature).
    \item That the gate-battery thresholds or the formulas in G2-1..3 have been independently audited by the arithmetic-geometry or theoretical physics communities (external verification is invited and is flagged as Phase 8.D task D-4, pending).
    \item That the Lean axioms (general-law declarations for the recurrences) are proven in-kernel (they are empirically validated to $n \approx 200$, declared as auditable assumptions, not `sorry` holes).
    \item That any observational signature of these candidates will be detected (PTA, Lyman--$\alpha$, or otherwise).
\end{enumerate}
These caveats are consistent with Parts I--VII's provenance-ledger discipline.
\end{remark}

\section{Ongoing and Deferred Items}

\subsection{Phase 8.D/D-3: Finalized LaTeX manuscript (this document)}

Completed 2026-07-14. Negative findings, methods, finalists, and discussion sections are now in place.

\subsection{Phase 8.D/D-5: Observatory targeting dossier}

Pending. Per-finalist PTA frequency bands and superradiance timescales, formatted for observational audiences, will be generated from \texttt{g2\_3\_superradiance\_bands.json} and \texttt{g2\_4\_obs\_screens.json}.

\subsection{Phase 8.D/D-4: External verification invitations}

Explicitly \textbf{not} executed in this session. Opening GitHub issues on external community repositories (arithmetic-geometry, PTA, axion dark-matter communities) is an outward-facing action that commits the programme's name externally and requires separate explicit user authorization. Flagged for Phase 8.E or later.

\section*{Acknowledgments}

The exact-arithmetic verification pipeline, Lean 4 formalization infrastructure, and monodromy integration code are built on top of the Mathlib community libraries and the mpmath numerical Python library. The gate-battery design follows the provenance-ledger discipline established in Parts I--VII of this series.

\begin{thebibliography}{99}

\bibitem{Cooper2012}
S.~Cooper,
``Sporadic sequences, modular forms and new series for $1/\pi$,''
\textit{Ramanujan J.} \textbf{29}, 163--183 (2012).

\bibitem{Beukers1985}
F.~Beukers,
``A note on the Irrationality of $\zeta(2)$ and $\zeta(3)$,''
\textit{Bull. London Math. Soc.} \textbf{11}, 268--272 (1979).

\bibitem{Beauville1982}
A.~Beauville,
``Variétés Kähleriennes dont la première classe de Chern est nulle,''
\textit{J. Differ. Geom.} \textbf{18}, 755--782 (1983).

\bibitem{Rogers2021}
K.~K.~Rogers and U.~L.~Peiris,
``Frequency-Dependent Strong Lensing in Fuzzy Dark Matter,''
\textit{Mon. Not. R. Astron. Soc.} \textbf{503}, 5956--5970 (2021).

\bibitem{Dolan2018}
S.~R.~Dolan,
``A separability ansatz for the gravitational Teukolsky equation in Kerr--Newman--de-Sitter spacetimes,''
\textit{Class. Quantum Grav.} \textbf{35}, 235007 (2018).

\end{thebibliography}

\end{document}
